% !TEX root = ../thesis.tex
\chapter*{TÓM TẮT KHÓA LUẬN}
\addcontentsline{toc}{chapter}{TÓM TẮT KHÓA LUẬN}

Sự bùng nổ của thông tin trực tuyến tại Việt Nam đã đặt ra nhu cầu cấp thiết về việc tóm tắt nhanh nội dung bài báo. Các hệ thống tóm tắt dựa trên một Mô hình ngôn ngữ lớn (LLM) đơn lẻ hiện nay vẫn cho ra kết quả không ổn định, trong khi kỹ thuật Mixture-of-Agents (MoA) --- kết hợp nhiều LLM thông qua kiến trúc đa tác nhân bao gồm các tác nhân đề xuất và tổng hợp --- chưa từng được nghiên cứu cho bài toán tóm tắt tin tức tiếng Việt. Bên cạnh đó, các hệ số đo lường phổ biến như ROUGE hay BERTScore chủ yếu đo độ trùng lặp với bài gốc, vốn kém hiệu quả trong việc phản ánh chất lượng khi mô hình được phép diễn giải lại và rút gọn ý tưởng.

Khóa luận này thực hiện ba đóng góp. Thứ nhất, giới thiệu hệ thống Fiber --- một tiện ích mở rộng trình duyệt có khả năng tóm tắt tin tức tiếng Việt theo thời gian thực, hỗ trợ các LLM đa nền tảng (OpenAI, Anthropic, Google). Thứ hai, áp dụng pipeline MoA cho bài toán tóm tắt tiếng Việt có neo ngữ cảnh gốc, bằng cách đưa trực tiếp bài viết nguồn vào tác nhân tổng hợp như một kết nối phần dư nhằm đảm bảo tính trung thực. Thứ ba, đề xuất khung đánh giá đa trục: Trục A (giữ lại nội dung: ROUGE, BERTScore), Trục B (chất lượng: đánh giá rubric và xếp hạng bởi LLM-Judge), và Trục C (đánh giá mù bởi con người), qua đó chỉ ra hạn chế của các hệ số overlap truyền thống.

Trên tập dữ liệu 154 bài báo, pipeline MoA cho kết quả tóm tắt vượt trội so với các mô hình đơn lẻ ($68,9\%$ tỷ lệ thắng). Các đánh giá định lượng và khảo sát độc lập trên 17 tình nguyện viên đều đồng thuận với kết quả này, đồng thời phản ánh rõ tác động của yếu tố chủ đề lên sở thích tiếp nhận thông tin của độc giả.
\vspace{0.4em}

Từ khoá: Tóm tắt tin tức tiếng Việt, mô hình ngôn ngữ lớn, Mixture-of-Agents, tiện ích mở rộng trình duyệt, khung đánh giá đa trục.
